\section{讨论}

路由实验回答第一个研究问题。DeepSeek Flash和GLM 5.3 Flash分别达到0.995和0.990的路由准确率，混合效用为0.403和0.419，后者恢复oracle的99.3\%。错路由主要发生在包含未来时态、但答案已经在文件中给出的确定问题。Qwen3-8B的路由准确率降至0.867，混合效用0.263，低于Always-Anchor的0.298。这个对比说明契约识别不是简单的时态匹配：较强基座能够利用查询与截止时点的关系，较小基座则容易把语义边界压缩为表面模式。路由实验因此支持H1在较强基座上成立，同时给出基座能力的边界。

DocFinQA和FinDVer回答第二个研究问题。DeepSeek Flash上，义务检索分支取得0.457精确匹配、平均10.69段，固定16段阅读器取得0.390、平均16段；GLM 5.3 Flash上分别为0.497对0.480，平均10.31段对16段。FinDVer上三个基座分别为0.701、0.753和0.649，均高于固定16段的0.660、0.743和0.636。四个具名RAG适配在DocFinQA上均低于义务检索分支，分支的答案覆盖率0.770和0.810也高于固定预算。规划器单独运行的成本更低，但在GLM 5.3 Flash上留下小缺口，置信度回退在24\%的问题上触发。实验支持H2：类型化义务与覆盖驱动取证在长文档上以更低证据成本取得更高准确率，并可迁移到FinDVer；收益来自证据策略，而不是阅读器或评分口径。

收入指引实验回答第三个研究问题，也给出最需要谨慎解释的结果。纯LLM锚点后验在三个基座上的准确率为0.445、0.454和0.427，决策损失为4.035、3.903和3.586，说明锚点包含非随机信号，但类别偏置和概率校准不足以直接承担低损失决策。融合后验把准确率提高到0.524、0.537和0.537，决策损失降到3.604、3.648和3.577，低于六个具名基线；在GLM 5.3 Flash和Qwen3-8B上，准确率也高于最强具名基线。纯RF诊断取得0.533、0.491、0.486和3.454，是当前最强的数值决策头。这个结果不支持“LLM锚点后验优于强RF”，但支持“语义状态与校准读出分工”：LLM负责路由、阅读和锚点状态，RF负责数值风险控制。H3应表述为融合读出相对具名基线降低决策损失，而不是LLM单独胜过RF。

消融实验进一步说明锚点层的作用与边界。去掉全部锚点后，LLM后验的宏F1从0.233、0.286、0.339降到0.213、0.211、0.272；取消硬锚点掩码后指标几乎不变，说明硬锚点在收入指引上很少触发；逐类去除锚点没有一致赢家，说明锚点类型的作用依赖基座和数据分布。输出形式对照显示，强制点答案的平均损失为4.693，固定三桶全集为4.000，损失感知门为3.454；风险收益来自集合与弃权，而不是更强的点分类器。误差分析显示，DocFinQA的错误主要来自表格块边界与算术，Qwen3-8B的错误主要来自解释性输出，AnchorEvolve的错误主要来自类别偏置，路由错误集中在语义边界。这些结果把改进方向限定在证据表示、锚点过滤和路由信号上，而不是盲目增加模型规模。

实验比较的是同基座适配，不是原系统的未修改端到端运行；标签来自XBRL核验而非人工金标准；混合负载由DocFinQA与收入指引样本构成。这些选择保证了阅读器、证据接口和评分口径可比，也意味着结论适用于给定的数据、基座和评价口径。它们不改变三个研究问题的检验结果，但界定了结果的外推范围。

\section{研究结论}

本文提出FinPlanRAG，把金融混合查询的回答形式选择表述为查询契约路由，并以义务检索和AnchorEvolve分别处理确定信息与前瞻性断言。路由器在检索与生成之前判断事实是否已在截止时点存在；义务检索围绕类型化义务形成闭合答案；AnchorEvolve以L0稳定关系、L1行业模式和L2实例锚点维护可更新、可追溯的前瞻性证据状态。校准读出在此基础上生成点答案、可信集或弃权，用于控制最终决策风险。

第一项贡献是查询契约路由。现有问答、断言验证和选择性预测分别固定输出接口，本文把“事实是否已在截止时点存在”置于检索与生成之前，用一次路由决定闭合答案或可校准判断。路由实验支持这一贡献：DeepSeek Flash和GLM 5.3 Flash分别达到0.995和0.990的路由准确率，混合效用接近oracle，并显著高于Always-RAG和Always-Anchor；Qwen3-8B的结果则说明路由能力依赖基座。契约路由的价值不只是提高一个分类指标，而是把混合负载的分析单位从数据集任务改为查询契约，使路由错误和分支错误可以分开诊断。

第二项贡献是义务检索与AnchorEvolve的两条证据分支。义务检索把确定问题绑定为类型化义务，在时点语料上按覆盖取证，由阅读器返回带证据编号的闭合答案；AnchorEvolve把前瞻性断言映射到有限结果空间，用三层锚点、五类资格判定和时点信念更新保留证据的作用域、可靠性和时效。DocFinQA和FinDVer实验支持义务检索的效率主张；收入指引实验和锚点消融支持AnchorEvolve的语义主张。特别是去除全部锚点后LLM后验宏F1下降，说明锚点不是装饰，而是改变了模型判断的中间状态。

第三项贡献是语义状态与校准读出的分工，以及统一的多基准评估。本文在三个数据集、三套基座和六个具名基线上同时检验路由、阅读、锚点推断、后验融合、校准选择和外部迁移，并把分类、校准、证据成本、决策损失和路由效用放在同一口径下报告。需要明确，纯RF是当前最强的数值决策头，融合后验在部分基座上取得更高普通准确率，但决策损失仍高于纯RF；因此本文不主张LLM锚点后验单独胜过强RF。本文主张的是：LLM在路由、阅读和语义状态构造上不可替代，校准读出负责把这些状态转换为低风险决策，二者应当分开评估。

本文的原创性贡献可以概括为三点。第一，把输出契约从固定接口变成查询层的路由变量，使确定问答与前瞻性裁决在同一时点约束下被统一处理。第二，把前瞻性证据从单一标签变成可审计的状态，使证据来源、作用域、资格理由和信念变化可以被逐条检查。第三，把语义判断与风险校准分开，使LLM的语义能力、统计模型的风险控制能力和路由器的契约识别能力可以分别检验、替换和组合。这三点不依赖特定基座或特定行业，构成了本文相对于现有金融RAG和断言验证工作的主要区别。

从金融实践看，框架把答案、证据和风险放在同一条审计链上。分析师对确定问题获得带证据编号的闭合答案，对前瞻性断言获得点判断、可信集或弃权，并可以查看锚点来源和资格理由。机构可以按自身风险偏好替换损失参数和校准后端，而不改变路由与证据状态的结构。生产里程碑、产能投放、监管审批和合同履约等任务与收入指引具有相似的时点结构，本文结果为这些场景提供可复用的接口。

本文的结论有明确边界。收入指引实验采用三桶结果空间，通用的时机与阶段状态仍需在更多前瞻性事件上评估；混合负载由DocFinQA与收入指引样本构成，真实用户查询的语义边界可能更复杂；标签来自XBRL核验而非人工金标准。后续工作将围绕更多前瞻性事件、更稳定的边界路由、更窄的可信集和更严格的人工复核展开。由于路由、锚点抽取和输出读出彼此分离，这些改进可以独立评估，而不需要改变FinPlanRAG的整体接口。
